\chapter{Grafos temporales: definiciones}
\label{cap:grafos-temporales}

\begin{intuicion}
Antes de hablar de \emph{matching} necesitamos un vocabulario
preciso. Un \emph{grafo dirigido temporal} es un conjunto de
vértices más un multiconjunto de aristas dirigidas con
\emph{timestamp}. ``Multi'' porque puede haber varias aristas entre
los mismos dos nodos (el mismo proceso lanzando otro varias veces);
``dirigido'' porque cada arista tiene origen y destino fijos (no son
simétricas); ``temporal'' porque el orden temporal de las aristas es
parte de la estructura, no un metadato. Este capítulo define cada
una de esas piezas con cuidado, distingue \emph{path} de \emph{walk}
(una trampa frecuente del threat hunting), y formaliza las
restricciones temporales que más tarde el matcher chequea
\emph{in-line}.
\end{intuicion}

\section{Grafos dirigidos: el caso base}

\begin{definition}[Grafo dirigido]
Un \emph{grafo dirigido} es un par \(G = (\V, \E)\) donde \(\V\) es
un conjunto finito de \emph{vértices} y \(\E \subseteq \V \times \V\)
es un conjunto finito de \emph{aristas dirigidas}. Una arista
\(e = (u, v)\) tiene \emph{origen} \(u\) y \emph{destino} \(v\). Si
\((u, v) \in \E\), \(v\) es un \emph{vecino saliente} de \(u\); \(u\)
es un \emph{vecino entrante} de \(v\).
\end{definition}

\noindent
Convención de tamaño: \(n = |\V|\) y \(m = |\E|\).

\begin{definition}[Vecindarios y grados]
Para \(v \in \V\):
\begin{itemize}
    \item \(N^{+}(v) = \{w : (v, w) \in \E\}\) es la \emph{vecindad
    saliente} de \(v\).
    \item \(N^{-}(v) = \{u : (u, v) \in \E\}\) es la \emph{vecindad
    entrante} de \(v\).
    \item \(\outdeg(v) = |N^{+}(v)|\) y \(\indeg(v) = |N^{-}(v)|\).
    \item \(\deg(v) = \outdeg(v) + \indeg(v)\).
\end{itemize}
\end{definition}

\section{Multigrafos}

En el caso base \(\E\) es un conjunto: existe a lo sumo una arista
entre cualesquiera dos vértices. En el threat hunting eso no alcanza:
el mismo proceso \(p\) puede crear el mismo archivo \(f\) cinco veces
en una hora; cada creación es un evento distinto.

\begin{definition}[Multigrafo dirigido]
Un \emph{multigrafo dirigido} es un par \(G = (\V, \E)\) donde \(\E\)
es un \emph{multiconjunto} sobre \(\V \times \V\). La
\emph{multiplicidad} \(\mu(u, v)\) es la cantidad de veces que el par
\((u, v)\) aparece en \(\E\).
\end{definition}

\begin{nota}
En la práctica, modelar GraphHunter como multigrafo es lo correcto:
distintos eventos entre las mismas dos entidades son distintos por
su \emph{timestamp} y su \emph{metadata}, aunque compartan
\((\text{origen}, \text{destino}, \text{tipo})\). El motor representa
cada evento como una arista distinta con su propio
\texttt{CompactRelation} (capítulo~\ref{cap:compact-relation}), no
como un contador sobre una arista compartida.
\end{nota}

\section{Etiquetas en vértices y aristas}

Hasta acá los vértices y aristas son indistinguibles entre sí. En el
threat hunting cada uno tiene un \emph{tipo}: un vértice puede ser un
\texttt{User}, un \texttt{Host}, una \texttt{IP}; una arista puede
ser una \texttt{Auth}, una \texttt{Connect}, un \texttt{Spawn}.

\begin{definition}[Grafo etiquetado]
Sean \(\Lc_V\) un alfabeto finito de \emph{etiquetas de vértice} y
\(\Lc_E\) un alfabeto finito de \emph{etiquetas de arista}. Un
\emph{grafo dirigido etiquetado} es una tupla
\((G, \mathrm{type}_V, \mathrm{type}_E)\) donde \(G = (\V, \E)\) es un
multigrafo dirigido,
\(\mathrm{type}_V \colon \V \to \Lc_V\) asigna una etiqueta a cada
vértice, y
\(\mathrm{type}_E \colon \E \to \Lc_E\) asigna una etiqueta a cada
arista.
\end{definition}

En GraphHunter \(\Lc_V\) es \texttt{EntityType} (\texttt{User},
\texttt{Host}, \texttt{Process}, \texttt{IP}, \texttt{File},
\texttt{Domain}, \texttt{Service}, \texttt{Resource},
\texttt{Group}\,---\,nueve tipos con \emph{tag} \([0, 9)\)) y
\(\Lc_E\) es \texttt{RelationType} (\texttt{Auth},
\texttt{Connect}, \texttt{Spawn}, \texttt{Read}, \texttt{Write},
\texttt{Beacon}, etc.).

\section{Grafos temporales}

Llegamos a la pieza que define al motor: el tiempo.

\begin{definition}[Grafo dirigido temporal]
Un \emph{grafo dirigido temporal etiquetado} es una tupla
\((\V, \E, \mathrm{type}_V, \mathrm{type}_E, t)\) donde
\((\V, \E, \mathrm{type}_V, \mathrm{type}_E)\) es un multigrafo
dirigido etiquetado y \(t \colon \E \to \T\) asigna un
\emph{timestamp} \(t(e)\) a cada arista. \(\T\) es el dominio
temporal: en GraphHunter, \(\T = \mathbb{Z}\) (segundos epoch).
\end{definition}

Cuando hablemos de \emph{una arista} en un grafo temporal usaremos la
notación \(e = (u, v, \tau)\), donde \(\tau = t(e)\). Igualdad de
aristas: dos aristas son distintas si difieren en
\(\mathrm{type}_E\), \(\tau\), o (en multigrafos) en cualquier
metadata adicional.

\begin{ejemplo}
Un grafo de seis vértices y siete aristas, todos con etiquetas y
\emph{timestamps}.

\begin{figure}[H]
\centering
\begin{tikzpicture}[
    every node/.style={font=\footnotesize},
    user/.style={draw=accent, thick, circle, minimum size=8mm,
                 fill=accentlight, inner sep=1pt},
    host/.style={draw=accent, thick, rectangle, rounded corners=1mm,
                 minimum size=8mm, fill=accentlight, inner sep=2pt},
    ip/.style={draw=accent, thick, diamond, minimum size=10mm,
               fill=accentlight, inner sep=1pt},
    edge/.style={->, >=Latex, accent, thick}
]
\node[user] (u1) {$u_1$};
\node[user, below=1.0cm of u1] (u2) {$u_2$};
\node[host, right=2.2cm of u1] (h1) {$h_1$};
\node[host, right=2.2cm of u2] (h2) {$h_2$};
\node[ip, right=2.2cm of h1] (ip1) {$\mathit{ip}_1$};
\node[ip, right=2.2cm of h2] (ip2) {$\mathit{ip}_2$};

\draw[edge] (u1) -- node[above]{\texttt{Auth}@1} (h1);
\draw[edge] (u2) -- node[above]{\texttt{Auth}@2} (h1);
\draw[edge] (u2) -- node[below]{\texttt{Auth}@5} (h2);
\draw[edge] (h1) -- node[above]{\texttt{Connect}@4} (ip1);
\draw[edge] (h1) -- node[pos=0.7,above,sloped]{\texttt{Connect}@7} (ip2);
\draw[edge] (h2) -- node[below]{\texttt{Connect}@6} (ip2);
\draw[edge] (h1) to[bend left=20] node[above]{\texttt{Connect}@9} (ip1);
\end{tikzpicture}
\caption{Grafo dirigido temporal etiquetado de ejemplo. \(\Lc_V =
\{\texttt{User}, \texttt{Host}, \texttt{IP}\}\), \(\Lc_E =
\{\texttt{Auth}, \texttt{Connect}\}\). Notar la \textbf{arista
múltiple} \(h_1 \to \mathit{ip}_1\): hay dos conexiones distintas, en
\(\tau = 4\) y \(\tau = 9\).}
\label{fig:gt-ejemplo}
\end{figure}
\end{ejemplo}

\section{Caminos: \emph{paths} vs \emph{walks}}

Es la distinción más importante del capítulo, y la trampa clásica.

\begin{definition}[Walk dirigido]
Un \emph{walk} dirigido de longitud \(k\) en \(G\) es una secuencia
de vértices \(v_0, v_1, \dots, v_k\) tal que para cada \(i \in
\{0, \dots, k-1\}\), existe una arista \((v_i, v_{i+1}, \tau_i) \in
\E\). Los vértices \emph{pueden repetirse}.
\end{definition}

\begin{definition}[Path simple]
Un \emph{path simple} (o sólo \emph{path}) es un \emph{walk} en el
que todos los vértices son distintos.
\end{definition}

\begin{advertencia}
En el threat hunting, \emph{walks} y \emph{paths} responden preguntas
distintas. ``¿Hay un usuario que se conectó a una IP y volvió a
conectarse a sí mismo?'' es una pregunta sobre \emph{walks} (el
usuario aparece en posiciones 0 y 2). ``¿Hay tres entidades distintas
encadenadas?'' es una pregunta sobre \emph{paths}. La DSL de
GraphHunter expresa esta elección con la cláusula \texttt{\{k=N\}}
(capítulo~\ref{cap:dsl-chain}): \(k\)-\emph{simplicidad} controla
cuántas veces puede repetirse un mismo vértice.
\end{advertencia}

\begin{definition}[\(k\)-simplicidad]
Un \emph{walk} \(v_0, \dots, v_k\) es \emph{\(k\)-simple} sii ningún
vértice aparece más de \(k\) veces en la secuencia. Caso particular:
\(1\)-simple = \emph{path simple}.
\end{definition}

\section{Restricciones temporales}

\subsection{Monotonicidad}

\begin{definition}[Walk temporalmente válido]
Un \emph{walk} \((v_0, e_0, v_1, e_1, \dots, v_k)\) ---donde
\(e_i = (v_i, v_{i+1}, \tau_i)\)--- es \emph{temporalmente válido} sii
\[
\tau_0 < \tau_1 < \dots < \tau_{k-1}.
\]
\end{definition}

Es decir: las aristas se atraviesan en orden cronológico estricto. Si
en el grafo de la figura~\ref{fig:gt-ejemplo} consideramos
\(u_2 \to h_2 \to \mathit{ip}_2\), las aristas tienen \(\tau\) igual a
5 y 6 ---walk válido. Pero \(u_2 \to h_1 \to \mathit{ip}_1\) tiene
\(\tau = 2\) y \(\tau = 4\) ---también válido. Y
\(h_1 \to \mathit{ip}_1\) (la primera, en \(\tau = 4\)), seguido de
\(h_1 \to \mathit{ip}_2\) (en \(\tau = 7\)), forma walk válido si
empezamos por una arista entrante a \(h_1\) anterior a \(\tau = 4\).

\begin{nota}
GraphHunter impone monotonicidad \emph{estricta} (\(<\), no \(\leq\)):
dos aristas con el mismo \emph{timestamp} no se ordenan entre sí.
Esto evita ambigüedad cuando dos eventos llegan en el mismo segundo
epoch (resolución típica del log de Sysmon). En la práctica, los
eventos genuinamente simultáneos son raros; cuando ocurren, el
matcher los descarta sin pérdida.
\end{nota}

\subsection{Ventana temporal}

\begin{definition}[Walk \(\Delta t\)-acotado]
Un walk temporalmente válido \((v_0, \dots, v_k)\) con timestamps
\((\tau_0, \dots, \tau_{k-1})\) está \emph{\(\Delta t\)-acotado} sii
\[
\tau_{k-1} - \tau_0 \leq \Delta t.
\]
\end{definition}

\(\Delta t\) es la \emph{ventana temporal} de la hipótesis. En la DSL
es la cláusula \texttt{WITHIN}; en el ejemplo Sysmon W0,
\(\Delta t = 30\)~minutos. Es la pieza que evita que el motor enumere
walks de tres saltos cuya primera arista es de hace tres meses.

\subsection{Composición: \emph{walk} de hipótesis}

\begin{definition}[Walk de hipótesis]
Dado un grafo temporal \(G\) y una hipótesis \(H\) con \emph{shape}
\emph{chain} de longitud \(k\), etiquetas
\((L_0^V, L_0^E, L_1^V, L_1^E, \dots, L_k^V)\), \(k\)-simplicidad
\(k_s\) y ventana \(\Delta t\), un \emph{walk de \(H\)} es una
secuencia
\((v_0, e_0, v_1, e_1, \dots, v_k)\) tal que:
\begin{enumerate}
    \item \(\mathrm{type}(v_i) = L_i^V\) para todo \(i\).
    \item \(\mathrm{type}(e_i) = L_i^E\) para todo \(i\).
    \item Es temporalmente válido: \(\tau_0 < \dots < \tau_{k-1}\).
    \item Está \(\Delta t\)-acotado: \(\tau_{k-1} - \tau_0 \leq
    \Delta t\).
    \item Es \(k_s\)-simple: ningún vértice se repite más de \(k_s\)
    veces.
    \item Cualquier predicado adicional sobre vértices o metadata se
    cumple.
\end{enumerate}
\end{definition}

El subgraph matching de GraphHunter es, esencialmente, enumerar
todos los walks de hipótesis para una \(H\) dada. El próximo
capítulo formaliza por qué eso es ---en general--- un problema duro,
y por qué en la práctica se vuelve tractable gracias a las
restricciones de tipo y temporales.

\section{Representación: \emph{adjacency list}}

Una decisión de implementación con consecuencias algorítmicas
profundas: GraphHunter no almacena \(\E\) como una matriz \(n \times
n\) (sería \(O(n^2)\) en memoria, intratable para \(n = 10^7\)).
Almacena \(\E\) como un \texttt{Vec<CompactRelation>} ordenado por
\((\texttt{source\_sid}, \texttt{timestamp})\), e indexa por
\texttt{source\_sid} para extraer rápidamente \(N^{+}(v)\).

\begin{enelcodigo}
\filepath{core/graph-engine/src/graph.rs}, función
\texttt{neighbors\_of(sid)}. Devuelve un \emph{slice} contiguo de
aristas salientes de \(v\), ya ordenadas por \texttt{timestamp}.
La búsqueda binaria sobre ese slice por la ventana temporal
\([t_{\mathrm{lo}}, t_{\mathrm{hi}})\) cuesta \(O(\log \outdeg(v))\)
y devuelve un sub-slice contiguo ---directamente iterable.
\end{enelcodigo}

Esa propiedad (\emph{contiguo}, \emph{ordenado},
\emph{binary-search-ready}) es la que el matcher explota. La Parte~V
del libro se apoya repetidamente en ella; el capítulo~\ref{cap:sistemas}
discute por qué la elección de un \texttt{Vec} contiguo en vez de un
\texttt{HashMap} de \texttt{Vec}s tiene impacto en \emph{cache}
performance.

\section{Resumen de notación}

\begin{center}
\renewcommand{\arraystretch}{1.2}
\begin{tabular}{p{3.0cm}p{9.5cm}}
\toprule
\textbf{Símbolo} & \textbf{Significado} \\
\midrule
\(G = (\V, \E, t)\) & grafo dirigido temporal etiquetado \\
\(n, m\) & \(|\V|\), \(|\E|\) \\
\((u, v, \tau)\) & arista temporal: origen \(u\), destino \(v\), timestamp \(\tau\) \\
\(N^{+}(v), N^{-}(v)\) & vecindad saliente, vecindad entrante \\
\(\outdeg(v), \indeg(v)\) & cardinalidad de cada vecindad \\
\(\Lc_V, \Lc_E\) & alfabetos de etiquetas (\texttt{EntityType}, \texttt{RelationType}) \\
\(\mathrm{type}(v), \mathrm{type}(e)\) & etiqueta del vértice / arista \\
\(k\) & longitud (en aristas) de un walk o de una hipótesis \\
\(k_s\) & cota de \(k\)-simplicidad \\
\(\Delta t\) & ventana temporal de la hipótesis \\
\bottomrule
\end{tabular}
\end{center}

\begin{ejercicios}
\begin{enumerate}
    \item \starred{} En el grafo de la figura~\ref{fig:gt-ejemplo},
    enumerá todos los walks de hipótesis con shape
    \texttt{User -[Auth]-> Host -[Connect]-> IP} sin restricción
    temporal. ¿Cuántos hay? Listalos como ternas \((u, h, \mathit{ip})\)
    junto con los timestamps.

    \item Repetí el ejercicio anterior con la restricción adicional
    \(\Delta t = 3\) (ventana de 3 unidades de tiempo). ¿Cuántos
    walks sobreviven?

    \item Construí un grafo temporal de exactamente 5 vértices y 6
    aristas que admita un \emph{walk} de longitud 4 pero ningún
    \emph{path simple} de longitud 4. (Pista: tiene que repetirse
    al menos un vértice.)

    \item \starred{} Definí formalmente la noción de
    ``\(\Delta t\)-conexidad temporal'': dos vértices \(u, v\) son
    \(\Delta t\)-conexos sii existe un walk temporalmente válido y
    \(\Delta t\)-acotado de \(u\) a \(v\). Mostrá que esta relación
    \emph{no} es de equivalencia en general (¿qué propiedad falla?).

    \item Si un grafo tiene \(n = 10^6\) vértices y \(m = 10^7\)
    aristas, ¿cuál es el grado promedio? ¿Cuánta memoria ocuparía la
    matriz de adyacencia? ¿Y la \emph{adjacency list}? ¿Por qué la
    diferencia importa para sistemas con \(n\) en los millones?

    \item Imaginá un caso de threat hunting donde la \(k\)-simplicidad
    \emph{importa}: una hipótesis donde permitir
    repeticiones (\(k_s = 2\)) cambia las respuestas legítimas.
    Después imaginá uno donde \(k_s = 1\) (path simple) es lo
    correcto. ¿Cuál es el riesgo de equivocarse?
\end{enumerate}
\end{ejercicios}
